\input{TD.sty}
\begin{document}
\Entete[oraux]{\theyears}

% Session 1
\clearpage
\begin{oral}[1]
  \Principal[amplitude]{Configuration lame d'air}
  \Probleme[referentiels]{La limite de Roche}
\end{oral}



\clearpage
\begin{oral}[1]
  \Principal*[parfait]{La clepsydre}
  \Probleme[interf]{Détermination du spectre d'une source lumineuse}
\end{oral}


% % Session 2
% \clearpage
% \begin{oral}[1]
%   \Principal[thermique]{Isolation thermique d'un bâtiment}
%   \Probleme[geometrique]{Détermination d’une distance focale}
% \end{oral}

% \clearpage
% \begin{oral}[1]
%   \Principal[solide]{Rotation d'un œuf dur}
%   \Probleme[thermique]{Combinaison de plongée}
% \end{oral}

% \clearpage
% \begin{oral}[2]
%   \Principal[electrostat]{Champ électrostatique d'un disque chargé}
%   \Ex[referentiels]{La limite de Roche}
% \end{oral}
% % %
% \clearpage
% \begin{oral}[2]
%   \Principal[amplitude]{Configuration lame d'air}
%   \Ex[induction]{Allumage d'une lampe par induction}
% \end{oral}
%
% \clearpage
% \begin{oral}[2]
%   \Principal[ondesEM]{Un exemple d'onde électromagnétique}
  % \Ex[parfait]{Lévitation d’une balle de ping-pong}
% \end{oral}

% \Ex[optique]{Interférences}
% \Ex[thermo]{Dissolution d’un grain de sucre}
% % \Ex{Sublimation d’une comète}
% \Ex[thermo]{Étude d'un détendeur adiabatique}
% \Ex[mecaflu]{Vidange d'un liquide visqueux}
% \Ex[meca]{Pendule de torsion}
% \Ex[electromag]{Supraconduction}
% \Ex[electromag]{Filtrage}
% \Ex[electromag]{Détection synchrone}
% \Ex[electromag]{Jonction entre semi-conducteurs}
% \Ex[meca]{Tir à la carabine}
% \Source{ondes.tex}
% \Ex{Transmission d’une onde acoustique}
\end{document}
